\subsection{Two-mark radial slope and phase rules}\label{c4:pair-radial}

We now derive the radial expression for the divided trace
\(\mathcal T(L,u_*)\) of \eqref{c4:confluent-trace}, retaining the slope
lattice and the phase of every root activity. This identifies the two
stationary operators on either side of the marked interval.
In particular, the central marked insertion
changes the slope by exactly \(3/2\), which will produce the decay in the
central radial region. Width in continuous variables is denoted by \(L\) (not the site block here). The deterministic vertices after passing to the conjugated cut and fusing the cluster are
\[
 B=b_H+2w,\qquad b_*= -2v,\qquad
 b_H=\sum_{i=1}^2(a_i h_i+d_i A h_i),\quad w=A^e v=C^e v,\quad e=0\ \text{or }1.
\]
Here \(a_i,d_i\) count unshifted and shifted colors in sector \(i\); the two marks have identical shift \(e\) in the application. Real positions are \(0,u_*\) with \(0<u_*<L\), label coefficients \(k_B=2,\ k_{b_*}=-2\) for the cocycle. Both angles may be taken as limiting angles 0; roots use the contours specified in the spread estimate. If the compensator lies angularly between parts of the approaching pure cluster, the two limiting orders with any cluster charge differ only by a deterministic unit phase: the direct power together with the shift power gives \((z-z')^{(b,b')}\) versus \((z'-z)^{(b,b')}\) times unit cocycles, analytically at distinct points since \(0<u_*<L\), while image products are unchanged. Thus the overall-unit convention for the fused field applies also then.

We also need stationary identities with no deterministic vertices, obtained from the same angular conjugation but using the additional character \(\exp(2\pi i(\sigma,P))\), \(\sigma=0\) or \(2v\). This commutes with the root actions in the conjugation. The normalized angular value in both cases is 1: before conjugation use just \(\widehat C\) (and the character), so only invariant momenta \(t\,4p_C,\ t\in\mathbb Z\), occur, with trivial character. In expanding the new cut place the character at the post-\(A\) momentum, then the unipotent factors (moving contours among external fields as in the angular argument). The case with deterministic insertions always starts with \(\sigma=0\).

Use the bilinear notations \(J(x,y)=x^t J y\), \(\mathcal D(x,y)=J(x,y)/2\) (field phase is \(\pi i[\mathcal D(b,P)-3k_b(p_C,P)]\)). Put
\[
 \begin{gathered}
 \alpha=\alpha_{10},\quad \beta=\alpha_{20},\quad
 \Pi=\Pi_A,\quad K=(A-I)^{-1}\text{ off the invariant line, }0\text{ on it},\\
 p=3\Pi v,\quad n(x)=(p,x),\quad s(x)=-(\alpha,Kx),\quad m_0(x)=-Kx+p\,s(x),\quad n_0=A m_0.
 \end{gathered}
\]
Here \(p^2=9/4,\ (\alpha,p)=-1,\ n(v)=n(w)=3/4,\ n(h_i)=0\), and \(p\) is primitive in the dual lattice. From Section~\ref{c4:pair-angular}, the \(D_1D_2\widehat A\) phase on a starting momentum \(m\), divided in the log by \(\pi i\), is \(F(m,m)+F_1(m)\), where
\[
 \begin{aligned}
 F(x,y)&=\tfrac12\sum_{r=\alpha,\beta}(r,Cx)\big((r,Cy)-J(r,Cy)\big),\\
 F_1(x)&=\tfrac12\sum_{i=1,2}(\alpha_{i0},Cx)(J(\alpha_{i0},h_i)-1)
                                  -\tfrac14(\alpha+\beta,A x).
 \end{aligned}
\]
To separate the momentum sum from the ordered field cocycles, use the
following forms. The form \(g\) records the quadratic phase at the
chosen representative \(m_0(d)\), and \(H\) will be the residual phase
of an ordered pair of charges. The corrections defining \(H^\#\) and
\(g^\#\) leave the total phase unchanged when \(n(d)=0\):
\begin{equation}
 \begin{aligned}
 U(x,y)&=F(x,y)+\mathcal D((I-A)x,Ay),\\
 g(x,y)&=F(m_0 x,m_0 y)+\mathcal D(x,n_0 y),\\
 g_{\sigma,1}(x)&=F_1(m_0 x)+2(\sigma,n_0 x),\qquad N(x,y)=(-Kx+\Pi x/2,y),\\
 H(x,y)&=g(x,y)+g(y,x)+\mathcal D(y,x)-N(x,y)+2n(x)s(y),\\
 H^\#(x,y)&=H(x,y)+H(x,\alpha)n(y)+H(y,\alpha)n(x)+H(\alpha,\alpha)n(x)n(y),
       \\ g^\#&=g+(H^\#-H)/2 .
 \end{aligned}                                                     \label{c4:PR-a}
\end{equation}

\paragraph{The phase identities needed for the transform.}
The forms in \eqref{c4:PR-a} separate the phase of a single root from
its residual pair phases. The following lemma gives the arithmetic
properties needed to make those residual phases trivial, and to identify
the momentum sum. We first use it to derive the radial rule; its full
coordinate verification closes this subsection.

\begin{lemma}[Finite phase reduction]\label{c4:finite-phase-reduction}
With the forms and charges just defined, let the columns of
\(\mathscr E=(I+\alpha n,\ h_1,h_2,Ah_1,Ah_2,(C-I)v)\)
consist of the images of a lattice basis under \(I+\alpha n\), followed
by the five displayed charges. Then
\par\noindent\textbullet\quad \(n\mathscr E=0,\ m_0\mathscr E\subset\mathcal L^*\) (columnwise),
\[
 U(p,p)=1,\quad F_1(p)=-1/2,\quad
              (U+U^t)(p,m_0\mathscr E)\equiv0\pmod{2}.
\]
\par\noindent\textbullet\quad \(H^\#(b,x),H^\#(x,b)\in2\mathbb Z\) for \(b\in\mathcal L\) and \(x\) in the integral span of \(\mathcal L,h_i,Ah_i,(C-I)v\). Define
\[
 r_e(x)=H^\#(w,x)-H^\#(x,v)+3(p_C,x)-2(v,n_0 x).
\]
Then \(r_0(I+\alpha n)\) is integral on the lattice, hence also \(r_e(I+\alpha n)\) for \(e=0,1\).
\end{lemma}

\paragraph{Poisson summation and local phases.} We derive the radial formula, using the earlier oscillator transformation \eqref{c4:Rad-1}--\eqref{c4:Rad-2} with \(R=A,\Pi,K\) as above for its mode functions. For each configuration of root charges and the indicated deterministic charges (or none), \(d=\sum b\), one requires \(n(d)=0\); then \(d\) is in the integral span of \(\mathscr E\). Thus all allowed pre-twist momenta \((I-A)m=d\) are exactly
\[
                          m=m_0(d)+t p,\qquad t\in\mathbb Z.
\]
The fields start from \(A m=m-d=K^*d+(t+s(d))p\), using \(K+K^*=-(I-\Pi)\). Before transforming, total momentum-dependent phase other than argument powers, in units \(\pi i\), is
\[
             U(m,m)+F_1(m)+2(\sigma,A m),
\]
since total \(k_b\) vanishes. By Lemma~\ref{c4:finite-phase-reduction},
using \(t^2-t\in2\mathbb Z\), its value modulo 2 is
\[
       g(d,d)+g_{\sigma,1}(d)+2t[1/4+(\sigma,p)] \pmod 2.
\]
The remaining cocycle phases are ordered pairs, value \(\mathcal D(b,a)-3 k_b(p_C,a)\) for angularly earlier \(a\) and later \(b\).

Apply the pair and self transforms exactly as for \eqref{c4:Rad-2}. They are valid on the same absolute-value-one eigenblocks (now also \(-1\)) with branch derivatives, by the derivation there. Their heat and order-independent linear argument corrections cancel off-invariant momentum powers/heat exactly since the post-twist state off-invariant has momentum \(K^*d\); the invariant Gaussian now uses \(t+s(d)\) with the extra \(t\)-phase above. Poisson summing gives initial dual label \(j_0\in1/4+(\sigma,p)+\mathbb Z\), with join scalar \(\sqrt{L/(2\pi p^2)}\) and phase \(-2 j_0 s(d)\) in the \(\pi i\) units. The real-position order terms give the same heat and angle rule: traverse rightward with \(j\) stepping by \(n(b)\) at each \(b\), apply
\begin{equation}
 \exp\!\left(-\frac1{2p^2}\int j(u)^2du\right),\qquad
                     \exp\!\left(i\phi_b n(b)\frac{j_{{\rm pre},b}+n(b)/2}{p^2}\right)
                    \quad\text{per angle}.
                                                                  \label{c4:PR-b}
\end{equation}
Here the base period contains \(0,u_*\) strictly inside when present. The image interactions are precisely \eqref{c4:Rad-2} and self weights are \(\exp((H_0[A]b,b)/2)\) as before. All root coefficient/integration conventions remain those of the angular expansion (factorwise exponential of zero modes, scaled fields with \(du/(2\pi)\)).

Besides these there are the pair sign phases of the transform, \(\operatorname{sgn}(u_a-u_b)N(a,b)\) for the same angular order \(a,b\). Thus when passed to radial order \(a\) then \(b\) we can always use
\[
                    \mathcal D(b,a)-N(a,b)-3 k_b(p_C,a)
\]
up to an overall deterministic unit. Indeed
\(N(x,y)+N(y,x)=(x,y)\); reversing from the other angular order
changes parity (when at least one charge is a root) by zero, by the
parity relation between \(J\) and the ordinary pairing, also valid for
shifted external charges by \(C\)-invariance. Self sign constants already cancel as in the earlier transform. Now use \(g(d,d)=g^\#(d,d)\), and write \(-2j_0s(d)=-2\sum j_{{\rm pre},b}s(b)+2\sum_{a<b} n(a)s(b)\) in radial order. This gives local phases
\begin{equation}
                g^\#(b,b)+g_{\sigma,1}(b)-2j_{{\rm pre},b}s(b)                 \label{c4:PR-c}
\end{equation}
in the same units, and residual ordered pairs \(H^\#(a,b)-3k_b(p_C,a)\).

\paragraph{The two stationary sectors.} For stationary traces all those residual pairs vanish modulo 2. For the two deterministic vertices with \(\sigma=0\), root-root and \(b_H\)-root parts vanish again. For a root \(b\) outside the interval \(0<u<u_*\) the pair effects from marks add to zero modulo 2 by the \((C-I)v\) integrality; inside they are
\[
       2[H^\#(w,b)-H^\#(b,v)+3(p_C,b)]\equiv
                      4(v,n_0 b)-2 r_e(\alpha)n(b) \pmod 2.
\]
Since \(n(B)=3/2=(2v,p)\), the slope sector inside is exactly that for the stationary \(\sigma=2v\) calculation, and the root factor phases there are \eqref{c4:PR-c} for that character, up to the diagonal-unit slope gauge with jump phase \(-2 r_e(\alpha)n(b)\). Thus they are stationary phases off and on the interval (with this fixed gauge), even through varied heavy profiles. Deterministic phases (including their mutual residual pair phase) only have modulus one.

\paragraph{The prefactor.} For the normalized prefactor put \(p_Q(r)=\mathfrak p(Q_L^r)\) as before. The characteristic polynomials give
\[
 Z_C=p_Q(2)/p_Q(8)^2,\qquad Z_A=p_Q(2)^3p_Q(3)^2/[p_Q(1)^2p_Q(6)^4].
\]
Under the scalar transform of these products, \(Z_C\) times its invariant momentum sum becomes
\(2\sqrt2\exp(-14\kappa/24-L/96)Z_C^r\Theta_C^r\), with \(\kappa=4\pi^2/L\), \(\Theta_C^r=\sum_{m\in\mathbb Z}\exp(-Lm^2/8)\); \(Z_A\) times the stated join scalar becomes \(2\sqrt2\exp(-14\kappa/24+L/48)Z_A^r\). Here radial oscillator products \(Z_R^r\) always use positive modes as in \eqref{c4:Rad-1}. Thus the full prefactor for \eqref{c4:PR-b}--\eqref{c4:PR-c} is
\begin{equation}
                         e^{L/32}\,\frac{Z_A^r(L)}{Z_C^r(L)\Theta_C^r(L)} .
                                                                  \label{c4:PR-d}
\end{equation}
With deterministic vertices \(B,-2v\), these rules give the radial
expression for \(\mathcal T(L,u_*)\), up to the fixed unit factor in the
confluence convention. In the stationary cases, with no deterministic
vertices, the radial series with this prefactor equals 1.

All identities use the formal-cut/absolute convergence and continuation bounds in the spread estimate: the angular quadratic bound suffices near small \(Q_L\) even with contour-to-mark commutations and then confluence as described above, and after momentum transformation the radial series has positive slope heat on the invariant line as required by \eqref{c4:Spread-3}, all residual phases unitary for real parameters. Thus the continuation argument in the final paragraph of
Section~\ref{c4:radial-coercivity}, using \eqref{c4:Spread-3} for the
two-mark bound, applies to the fused vertex and the two stationary
identities for every real $L>0$.

In passing to long radial links one can now remove \(b_H\) from the central vertex, extracting the real self scale \(\mathcal S\) of \eqref{c4:heavy-self-scale}, unit phases of bounded modulus even on varying slopes, interactions with mark/images tending to 1 in the successive limits at fixed site sizes (except for the counted self scale), and multipliers on root activities as in the spread estimate. Indeed \((A+A^{-1})h_i=h_i,\ n(b_H)=0\), and \((A^j h_i,v)=(A^j h_i,h_{i'})\) is exactly the displayed six-offset \(H'\). Thus all heavy self powers including cross terms with \(2w\) are the ones used there; \(H_0[A]\) contributions only among the two marked insertions themselves are bounded independent of site counts. Off the central inserting link the remaining root phases are literally the stationary ones on their slope sectors up to the unit gauge noted above. Independent of other phases, the entrance-to-exit slope across the central mark steps upward by \(3/2\).

\paragraph{Coordinate verification of the phase identities.}
\begin{proof}[Proof of Lemma~\ref{c4:finite-phase-reduction}]
Use the lattice basis \(\alpha_{1,0:7},\alpha_{2,0:5}\); the omitted even (odd) generator of sector 2 is minus the sum of the other even (odd) generators there and the odd (even) generators of sector 1. Thus \(C\) just shifts with this substitution, and the forms use the leading submatrices \(G,J\) of the circulants above. In these coordinates one has
\[
 \begin{split}
 3h_1^t&=(2,1,-1,-2,-1,1,2,1;\ 0,0,0,0,0,0),\\
 3h_2^t&=(-1,-2,-1,-2,-1,-2,-1,-2;\ 0,0,-3,-3,-3,0),\\
 24v^t&=(1,-13,-5,-19,-11,-1,7,-7;\ -6,-6,-12,-12,-18,6),\\
 n&=(-1,0,-1,1,-1,1,0,1;\ -1,0,-1,1,-1,1).
 \end{split}
\]
By \(A=(I-\alpha(\alpha,\cdot)-\beta(\beta,\cdot))C\), its characteristic polynomial is \((z-1)S(z)\), \(S=(z+1)(z^2-z+1)^4(z^2+z+1)^2\); \((A^6-I)^3=0\). For substitution one can use
\[
\begin{gathered}
 c_{0:13}=[1,-1,3,-1,4,0,3,3,0,4,-1,3,-1,1],\\
 \Pi=\sum c_j A^j/18,\qquad K=(117\Pi-\sum_j c_j\sum_{l<j,\ l\ge0}A^l)/18 .
\end{gathered}
\]
(The \(c_j\) are coefficients of \(S\).) For clarity, multiplying \(G(I+p(\alpha,\cdot))(-K)\mathscr E\) gives rows
\begingroup\setlength{\arraycolsep}{2pt}\scriptsize
\[\begin{array}{rrrrrrrrrrrrrrrrrrr}
0 & 0 & 0 & 0 & 0 & 0 & 0 & 0 & 0 & 0 & 0 & 0 & 0 & 0 & 0 & 0 & 0 & 0 & 0\\
0 & 1 & -1 & 2 & -1 & 1 & 1 & 1 & -1 & 0 & -1 & -1 & -1 & 1 & 1 & 0 & 1 & 0 & 0\\
0 & 1 & 1 & 0 & 1 & 0 & 0 & 1 & 0 & 0 & 0 & 0 & -2 & 0 & 0 & 0 & 1 & 0 & 0\\
0 & -2 & -2 & 2 & -4 & 1 & 0 & -1 & -2 & 2 & 0 & 0 & 2 & 0 & 0 & 0 & -1 & 2 & 1\\
0 & 2 & 1 & 1 & 2 & 0 & 1 & 1 & 0 & -2 & 0 & 0 & -2 & 0 & 0 & -2 & 1 & -2 & 0\\
0 & -2 & -1 & 2 & -3 & 2 & 0 & 0 & -2 & 2 & -2 & 0 & 2 & 0 & 0 & 0 & -1 & 0 & 0\\
0 & -1 & -1 & 1 & -1 & 0 & 1 & -1 & -1 & 0 & 1 & -1 & 1 & 1 & 0 & 0 & -1 & 0 & -1\\
0 & -1 & -1 & 2 & -2 & 2 & 1 & 1 & -2 & 0 & -2 & 0 & 0 & 0 & 1 & 0 & 0 & 0 & 0\\
0 & 2 & 1 & -2 & 4 & -2 & 0 & 1 & 2 & -2 & 1 & 0 & -2 & 0 & 0 & 0 & 1 & -1 & 0\\
0 & 0 & 0 & -2 & 0 & 0 & -2 & 0 & 1 & 1 & 0 & 1 & 0 & 0 & 0 & 1 & 0 & 1 & 0\\
0 & 2 & 1 & -2 & 2 & -2 & 0 & -1 & 2 & -1 & 2 & 0 & -1 & 0 & 0 & 0 & 1 & 0 & 0\\
0 & 0 & 1 & 0 & 0 & 0 & 0 & 1 & 0 & 0 & -1 & 0 & 0 & -1 & 0 & 0 & 1 & 0 & 1\\
0 & 0 & 1 & -2 & 2 & -2 & -2 & -1 & 2 & 0 & 2 & 1 & 0 & 0 & -2 & 0 & -1 & 0 & 0\\
0 & 0 & -1 & 0 & 0 & 0 & 0 & -1 & 0 & 0 & 0 & 0 & 1 & 0 & 0 & 0 & -1 & 0 & 0\\
\end{array}\]
\endgroup
and \((U+U^t)(p,\cdot)=(-2,0,-2,2,-2,2,0,2;\ 0,2,-2,2,-2,0)\). Likewise \eqref{c4:PR-a} gives \(H^\#/2\) with rows
\begingroup\setlength{\arraycolsep}{2pt}\scriptsize
\[\begin{array}{rrrrrrrrrrrrrr}
0 & 0 & 0 & -1 & 1 & -1 & 0 & 0 & 1 & 0 & 1 & 0 & 0 & 0\\
0 & 1 & -1 & -1 & -1 & -1 & -1 & 1 & -1 & 0 & 1 & 2 & 0 & 2\\
0 & 1 & 0 & 0 & 0 & -1 & 0 & 0 & 0 & 0 & 1 & 0 & 0 & 1\\
0 & -3 & 0 & 4 & -4 & 3 & 1 & -1 & -3 & 1 & -3 & -1 & 2 & -1\\
0 & 1 & 0 & -1 & 2 & -2 & -1 & 0 & 0 & 0 & 3 & 1 & 0 & 1\\
0 & -3 & -1 & 4 & -3 & 4 & 2 & -1 & -3 & 0 & -3 & -1 & 2 & -1\\
0 & -2 & 0 & 1 & -1 & 2 & 1 & -1 & -1 & 1 & -1 & -1 & 2 & -1\\
0 & -1 & -1 & 1 & -2 & 1 & 1 & 0 & -2 & 0 & 0 & 0 & 1 & 1\\
0 & 1 & 0 & -3 & 2 & -3 & -1 & 0 & 2 & 0 & 3 & 1 & 0 & 1\\
0 & 0 & 0 & -1 & 1 & -1 & 0 & 0 & 1 & -1 & 0 & 0 & 0 & 0\\
0 & 3 & 1 & -4 & 4 & -3 & -2 & 1 & 3 & -1 & 2 & 1 & -2 & 1\\
0 & 2 & 0 & 0 & 0 & 0 & 0 & 1 & 0 & -1 & -1 & 0 & -2 & 1\\
0 & 1 & 2 & -1 & 3 & -1 & 0 & 0 & 3 & 0 & 0 & -2 & -2 & -2\\
0 & 1 & 0 & 0 & 0 & 0 & -1 & 1 & 0 & 0 & 0 & 1 & -2 & 0\\
\end{array}\]
\endgroup
Multiplying this matrix in either argument by the coordinate columns of
\(h_i,Ah_i,(C-I)v\) gives integral columns and rows. This proves the
evenness of \(H^\#\) on the extended span in both argument positions.
Finally,
\[
 r_0(I+\alpha n)=(0,0,0,0,2,2,2,0;\ 0,0,0,-2,-2,-2).
\]
The difference \(r_1-r_0=H^\#((C-I)v,\cdot)\) is even on the lattice,
so the same integrality holds for both values of \(e\). These
computations verify all the stated identities.
\end{proof}
